
\documentclass[12pt,a4paper]{article}

% =========================================================
% PAQUETES
% =========================================================
\usepackage[utf8]{inputenc}
\usepackage[T1]{fontenc}
\usepackage[spanish,es-nodecimaldot]{babel}
\usepackage{lmodern}
\usepackage{geometry}
\usepackage{graphicx}
\usepackage{booktabs}
\usepackage{array}
\usepackage{longtable}
\usepackage{float}
\usepackage{amsmath,amssymb}
\usepackage{enumitem}
\usepackage{xcolor}
\usepackage{titlesec}
\usepackage{fancyhdr}
\usepackage{hyperref}
\usepackage{setspace}
\usepackage{lastpage}
\usepackage{caption}
\usepackage{subcaption}
\usepackage{tikz}
\usepackage{tabularx}
\usepackage{pdflscape}

% =========================================================
% CONFIGURACIÓN GENERAL
% =========================================================
\geometry{
    top=2.5cm,
    bottom=2.5cm,
    left=2.7cm,
    right=2.7cm
}

\setstretch{1.15}
\setlength{\parindent}{0pt}
\setlength{\parskip}{0.6em}

% Color institucional aproximado para títulos.
% Puede modificarse si se desea emplear otro tono.
\definecolor{UBUblue}{RGB}{0,82,147}
\definecolor{lightgray}{RGB}{245,245,245}
\definecolor{UBUlight}{RGB}{235,243,249}
\definecolor{UBUdark}{RGB}{34,45,58}

% Carpeta donde se encuentra el logotipo institucional
\graphicspath{{logo/}}

\hypersetup{
    colorlinks=true,
    linkcolor=UBUblue,
    urlcolor=UBUblue,
    citecolor=UBUblue,
    pdftitle={Trabajo - Programación Orientada a Objetos},
    pdfauthor={Universidad de Burgos}
}

% =========================================================
% DATOS DEL TRABAJO
% MODIFICAR ESTA SECCIÓN
% =========================================================

\newcommand{\universidad}{Universidad de Burgos}
\newcommand{\grado}{Grado en Tecnologías Digitales para la Empresa}
\newcommand{\asignatura}{Programación Modular y Orientada a Objetos}
\newcommand{\titulotrabajo}{Refactorización de una planta industrial mediante POO}

\newcommand{\nombrealumno}{Pendiente de completar}
\newcommand{\dni}{Pendiente de completar}
\newcommand{\correo}{Pendiente de completar}
\newcommand{\grupo}{Pendiente de completar}
\newcommand{\profesor}{Pendiente de completar}
\newcommand{\cursoacademico}{2026--2027}
\newcommand{\fechaentrega}{\today}

% =========================================================
% FORMATO DE TÍTULOS
% =========================================================

\titleformat{\section}
  {\Large\bfseries\color{UBUblue}}
  {\thesection.}{0.6em}{}

\titleformat{\subsection}
  {\large\bfseries\color{UBUblue}}
  {\thesubsection.}{0.6em}{}

\titleformat{\subsubsection}
  {\normalsize\bfseries}
  {\thesubsubsection.}{0.6em}{}

% =========================================================
% CABECERA Y PIE
% =========================================================

\pagestyle{fancy}
\fancyhf{}

\fancyhead[L]{\small \asignatura}
\fancyhead[R]{\small \grado}

\fancyfoot[L]{\small \universidad}
\fancyfoot[C]{\small \nombrealumno}
\fancyfoot[R]{\small Página \thepage\ de \pageref{LastPage}}

\setlength{\headheight}{15pt}
\renewcommand{\headrulewidth}{0.4pt}
\renewcommand{\footrulewidth}{0.4pt}

% =========================================================
% DOCUMENTO
% =========================================================

\begin{document}

% =========================================================
% PORTADA
% =========================================================

\begin{titlepage}
    \thispagestyle{empty}

    % Elementos gráficos de fondo
    \begin{tikzpicture}[remember picture,overlay]
        % Banda institucional izquierda
        \fill[UBUblue]
            (current page.north west)
            rectangle
            ([xshift=1.15cm]current page.south west);

        % Línea superior de acento
        \fill[UBUblue]
            ([xshift=1.15cm,yshift=-0.8cm]current page.north west)
            rectangle
            ([yshift=-0.95cm]current page.north east);

        % Bloque inferior suave
        \fill[UBUlight]
            ([xshift=1.15cm,yshift=4.0cm]current page.south west)
            rectangle
            (current page.south east);
    \end{tikzpicture}

    \hspace*{0.8cm}
    \begin{minipage}[t]{0.89\textwidth}

        \vspace*{0.2cm}

        % -------------------------------------------------
        % LOGO UBU
        % El archivo debe estar en: logo/logo_ubu.png
        % -------------------------------------------------
        \IfFileExists{logo/logo_ubu.png}{%
            \includegraphics[width=0.64\textwidth]{logo/logo_ubu.png}%
        }{%
            \fbox{%
                \parbox[c][2.8cm][c]{0.62\textwidth}{%
                    \centering
                    \textbf{LOGO UNIVERSIDAD DE BURGOS}\\[0.25cm]
                    Añadir \texttt{logo/logo\_ubu.png}%
                }%
            }%
        }

        \vspace{1.25cm}

        {\small\bfseries\color{UBUblue}\MakeUppercase{\grado}}\\[0.22cm]
        {\large\color{UBUdark}\universidad}

        \vspace{1.55cm}

        {\small\bfseries\color{UBUblue}\MakeUppercase{Asignatura}}\\[0.18cm]
        {\LARGE\bfseries\color{UBUdark}\asignatura}

        \vspace{0.75cm}
        {\color{UBUblue}\rule{0.22\textwidth}{3pt}}

        \vspace{1.0cm}

        {\small\bfseries\color{UBUblue}\MakeUppercase{Trabajo / práctica}}\\[0.30cm]
        {\Huge\bfseries\color{UBUdark}\titulotrabajo}

        \vfill

        % -------------------------------------------------
        % BLOQUE DE DATOS DEL ESTUDIANTE
        % -------------------------------------------------
        \noindent
        \colorbox{UBUlight}{%
            \parbox{0.96\linewidth}{%
                \vspace{0.35cm}
                \begin{tabularx}{\linewidth}{>{\bfseries\color{UBUdark}}p{4.0cm} X}
                    Estudiante & \nombrealumno \\[0.20cm]
                    DNI / NIE & \dni \\[0.20cm]
                    Correo electrónico & \correo \\[0.20cm]
                    Grupo & \grupo \\[0.20cm]
                    Profesor/a & \profesor \\[0.20cm]
                    Curso académico & \cursoacademico \\[0.20cm]
                    Fecha de entrega & \fechaentrega \\
                \end{tabularx}
                \vspace{0.15cm}
            }%
        }

        \vspace{0.85cm}

        {\small\color{UBUdark}
        Campus de Miranda de Ebro\hfill
        \textbf{Universidad de Burgos}}

        \vspace*{0.65cm}

    \end{minipage}
\end{titlepage}

% =========================================================
% ÍNDICE
% =========================================================

\tableofcontents
\newpage

% =========================================================
% 1. INTRODUCCIÓN
% =========================================================

\section{Introducción}

Esta práctica refactoriza una simulación funcional de una planta industrial,
originalmente organizada de forma principalmente procedural, para representar
sus componentes mediante objetos con responsabilidades diferenciadas. La planta
simulada en MuJoCo contiene cuatro cintas que forman un circuito, una quinta
cinta que conduce al horno, piezas de distintas geometrías, cámaras de
inspección y un manipulador Franka Emika Panda. OpenCV se utiliza para procesar
las imágenes capturadas desde las cámaras.

La solución mantiene el simulador y las rutinas físicas proporcionadas, pero
añade una capa orientada a objetos para coordinar la planta, encapsular el
estado de las cintas y del horno, representar las piezas y asociar cada cámara
con su cinta. La cinemática inversa (IK) del Panda se reutiliza como componente
existente; no se desarrolla ni se sustituye su algoritmo.

El punto de entrada es \texttt{main.py}. La lógica orientada a objetos reside
principalmente en \texttt{src/}, mientras que \texttt{planta\_4\_cintas.py}
conserva funciones auxiliares de simulación y configuración procedentes de la
base del proyecto. La aplicación está escrita en Python y combina MuJoCo,
NumPy y OpenCV.

% =========================================================
% 2. OBJETIVOS
% =========================================================

\section{Objetivos}

\subsection{Objetivo general}

Reorganizar la simulación de la planta industrial mediante clases con
responsabilidades claras, manteniendo su funcionamiento y preparando una
arquitectura modular que pueda ampliarse en la Práctica II.

\subsection{Objetivos específicos}

\begin{enumerate}
    \item Encapsular el estado y las operaciones de cada cinta transportadora.
    \item Asociar cámaras de inspección con sus cintas y procesar imágenes con OpenCV.
    \item Representar las piezas y sus datos relevantes para la simulación.
    \item Gestionar la recepción de piezas, la temperatura y la representación visual del horno.
    \item Coordinar los componentes desde \texttt{PlantaIndustrial} y conservar la integración con MuJoCo.
    \item Reutilizar el controlador IK del robot Panda sin reimplementar su cinemática.
\end{enumerate}

% =========================================================
% 3. ANÁLISIS DEL PROBLEMA
% =========================================================

\section{Análisis del problema}

El código de partida ya resolvía la simulación física y el transporte de piezas,
pero distribuía responsabilidades entre funciones, configuración global y
estructuras compartidas. La práctica pide trasladar el estado propio de cada
elemento a objetos, mantener separada la entrada del programa y lograr que la
detección de una cámara pueda influir en la cinta correspondiente. También
requiere que la temperatura del horno se refleje en su geometría de MuJoCo.

\subsection{Requisitos funcionales}

\begin{enumerate}
    \item Ejecutar la planta y visualizar la escena mediante MuJoCo.
    \item Gestionar de forma independiente cinco cintas, con controles de marcha y parada.
    \item Capturar imágenes de cinco cámaras y detectar presencia y forma de piezas.
    \item Actualizar la variable booleana \texttt{objeto\_detectado} de cada cámara y detener la cinta asociada durante la inspección.
    \item Registrar las piezas que llegan al horno, actualizar su temperatura y procesarlas al alcanzar el umbral definido.
    \item Cambiar el color de la geometría del horno según su temperatura.
    \item Mantener disponible el controlador Panda y su IK proporcionados.
\end{enumerate}

\subsection{Restricciones y consideraciones}

La solución debe continuar utilizando MuJoCo para la simulación y OpenCV para
la visualización y el procesamiento de cámaras. La IK del robot es código
proporcionado y queda fuera del desarrollo matemático de esta práctica. Las
funciones auxiliares y constantes físicas siguen centralizadas en
\texttt{planta\_4\_cintas.py}; por tanto, la refactorización reduce el uso de
estado global para los componentes modelados, pero no elimina toda la
configuración procedural heredada. El enunciado requiere que el informe final
incluya capturas y un enlace al repositorio GitHub; ninguno de esos elementos
externos está presente en los archivos disponibles para completar esta memoria.

% =========================================================
% 4. DISEÑO ORIENTADO A OBJETOS
% =========================================================

\section{Diseño orientado a objetos}

La clase coordinadora crea y mantiene las instancias de cintas, cámaras,
piezas, horno y robot. Cada cámara conserva una referencia a su cinta; la
planta sincroniza el estado lógico de las piezas con los cuerpos de MuJoCo y
actualiza las fuerzas, actuadores, cámaras y controles. Las rutinas físicas
auxiliares siguen delegadas en el módulo preexistente.

\subsection{Identificación de clases}

\begin{table}[H]
    \centering
    \caption{Clases principales de la aplicación.}
    \label{tab:clases}
    \begin{tabularx}{\textwidth}{p{3.2cm}X X}
        \toprule
        \textbf{Clase} & \textbf{Responsabilidad} & \textbf{Elementos principales} \\
        \midrule
        \texttt{PlantaIndustrial} & Coordina la simulación, los componentes y la interfaz. & \texttt{inicializar()}, \texttt{actualizar()}, \texttt{ejecutar()}, controles de teclado. \\
        \texttt{CintaTransportadora} & Mantiene el estado, la dirección y los objetos asociados a una cinta. & \texttt{iniciar()}, \texttt{detener()}, \texttt{conmutar()}, \texttt{esta\_activa()}. \\
        \texttt{CamaraInspeccion} & Captura imágenes, clasifica contornos y comunica detecciones. & \texttt{capturar()}, \texttt{detectar\_forma()}, \texttt{procesar()}, \texttt{objeto\_detectado}. \\
        \texttt{ObjetoIndustrial} & Representa una pieza, su forma, segmento, posición y orientación. & \texttt{actualizar()}, \texttt{cambiar\_segmento()}, \texttt{obtener\_pose()}. \\
        \texttt{Caja}, \texttt{Cilindro} & Especializan la construcción de piezas según su forma. & Heredan de \texttt{ObjetoIndustrial}. \\
        \texttt{Horno} & Registra y procesa piezas y gestiona temperatura y color. & \texttt{registrar\_pieza()}, \texttt{actualizar()}, \texttt{actualizar\_color\_mujoco()}. \\
        \texttt{RobotPanda} & Adapta el controlador IK existente al coordinador de la planta. & Inicialización, control manual/automático y sincronización de pose. \\
        \bottomrule
    \end{tabularx}
\end{table}

\subsection{Diagrama de clases UML}

\begin{landscape}
\begin{figure}[H]
    \centering
    \begin{tikzpicture}[
        clase/.style={draw, rounded corners, align=left, text width=4.0cm,
                      inner sep=7pt, font=\small},
        relacion/.style={-stealth, thick},
        escala/.style={font=\small}
    ]
        \node[clase] (planta) at (0,0) {\textbf{PlantaIndustrial}\\
            cintas, cámaras, objetos\\horno, robot\\
            \rule{3.7cm}{0.2pt}\\
            inicializar()\\actualizar()\\ejecutar()};
        \node[clase] (camara) at (-7.0,-4.5) {\textbf{CamaraInspeccion}\\
            nombre, cinta\\objeto\_detectado\\forma\_detectada\\
            \rule{3.7cm}{0.2pt}\\
            capturar()\\detectar\_forma()};
        \node[clase] (cinta) at (0,-4.5) {\textbf{CintaTransportadora}\\
            nombre, velocidad\\dirección, activa, objetos\\
            \rule{3.7cm}{0.2pt}\\
            iniciar()\\detener()\\conmutar()};
        \node[clase] (horno) at (7.0,-4.5) {\textbf{Horno}\\
            temperatura, piezas\\cola, geom\_id\\
            \rule{3.7cm}{0.2pt}\\
            registrar\_pieza()\\actualizar\_color()};
        \node[clase] (objeto) at (0,-9.0) {\textbf{ObjetoIndustrial}\\
            id, forma, segmento\\posición, orientación\\
            \rule{3.7cm}{0.2pt}\\
            actualizar()\\obtener\_pose()};
        \node[clase] (robot) at (7.0,0) {\textbf{RobotPanda}\\
            pose, modelo, estado\\
            \rule{3.7cm}{0.2pt}\\
            controlar\_manual()\\actualizar\_ciclo()};
        \draw[relacion] (planta) -- node[escala, left]{gestiona} (camara);
        \draw[relacion] (planta) -- node[escala, right]{gestiona} (cinta);
        \draw[relacion] (planta) -- node[escala, right]{gestiona} (horno);
        \draw[relacion] (planta) -- node[escala, above]{gestiona} (robot);
        \draw[relacion] (camara) -- node[escala, above]{asociada a} (cinta);
        \draw[relacion] (cinta) -- node[escala, right]{transporta} (objeto);
        \draw[relacion] (horno) -- node[escala, right]{procesa} (objeto);
    \end{tikzpicture}
    \caption{Diagrama de clases UML de la aplicación.}
    \label{fig:uml}
\end{figure}
\end{landscape}

\subsection{Relaciones entre clases}

\begin{itemize}
    \item \textbf{Composición y agregación:} \texttt{PlantaIndustrial} mantiene colecciones de cintas, cámaras y piezas, además del horno y el adaptador del robot. Las cámaras reciben una cinta asociada.
    \item \textbf{Asociación:} las cintas conservan las piezas que circulan por ellas; el horno mantiene una cola de piezas registradas.
    \item \textbf{Herencia:} \texttt{Caja} y \texttt{Cilindro} derivan de \texttt{ObjetoIndustrial} y fijan la forma inicial.
    \item \textbf{Dependencia:} \texttt{RobotPanda} delega el cálculo y control IK en el módulo \texttt{robot.py}; \texttt{PlantaIndustrial} utiliza las funciones físicas de \texttt{planta\_4\_cintas.py}.
\end{itemize}

% =========================================================
% 5. PRINCIPIOS DE POO APLICADOS
% =========================================================

\section{Principios de programación orientada a objetos aplicados}

\subsection{Encapsulación}

El estado de cada cinta (activa, velocidad, dirección y objetos) reside en su
instancia y se modifica mediante métodos de marcha, parada y conmutación. La
cámara mantiene sus resultados de inspección y referencia a la cinta asociada.
El horno agrupa la temperatura, sus límites, el registro de piezas y el
identificador de la geometría visual. La encapsulación es principalmente
convencional en Python: muchos atributos son accesibles directamente; la cola
del horno y algunos detalles de coordinación se indican como internos mediante
el prefijo \texttt{\_}.

\subsection{Abstracción}

Las clases expresan conceptos de la planta sin exponer al usuario de la
aplicación los detalles de cada operación: por ejemplo, \texttt{detener()} no
necesita conocer los actuadores de MuJoCo, y \texttt{registrar\_pieza()} agrupa
las actualizaciones del horno. La planta coordinadora oculta la secuencia de
inicialización, actualización física y cierre de ventanas.

\subsection{Herencia}

\texttt{Caja} y \texttt{Cilindro} heredan de \texttt{ObjetoIndustrial}. En esta
implementación, las subclases especializan sobre todo el constructor fijando la
forma, y reutilizan la identidad, segmento, posición, orientación y operaciones
comunes definidas en la clase base.

\subsection{Polimorfismo}

No se define una jerarquía con sobrescritura de comportamiento ni clases
abstractas. La herencia se utiliza para compartir la inicialización y el estado
común de las piezas; por ello, no se atribuye a la solución un uso significativo
de polimorfismo.

% =========================================================
% 6. IMPLEMENTACIÓN
% =========================================================

\section{Implementación}

La aplicación construye el modelo y los datos de MuJoCo, enlaza las clases con
los cuerpos y actuadores de la escena, ejecuta el ciclo de simulación y libera
los recursos gráficos al terminar. El ciclo de control incluye el movimiento de
las piezas, la interacción con el robot, la actualización del horno y la
captura periódica de cámaras.

\subsection{Estructura del proyecto}

\begin{itemize}
    \item \texttt{main.py}: crea \texttt{PlantaIndustrial} y llama a \texttt{ejecutar()}.
    \item \texttt{src/planta\_industrial.py}: coordinación del simulador y de sus componentes.
    \item \texttt{src/cinta\_transportadora.py}, \texttt{src/camara\_inspeccion.py}, \texttt{src/objeto\_industrial.py} y \texttt{src/horno.py}: clases de dominio.
    \item \texttt{src/robot.py}: adaptador orientado a objetos del controlador del Panda.
    \item \texttt{src/planta.xml}: geometría, cuerpos y cámaras MuJoCo de la planta.
    \item \texttt{planta\_4\_cintas.py}: constantes y funciones auxiliares físicas conservadas del simulador inicial.
    \item \texttt{robot.py}: controlador IK reutilizado.
    \item \texttt{requirements.txt}: dependencias Python; \texttt{README.md}: instalación, ejecución y controles.
    \item \texttt{vendor/mujoco\_menagerie/}: modelo y recursos del Franka Panda.
\end{itemize}

\subsection{Clases y métodos principales}

\texttt{PlantaIndustrial} configura las cinco cintas y sus cámaras, crea las
piezas, inicializa MuJoCo y el robot, y actualiza los componentes durante la
ejecución. Además gestiona las teclas: \texttt{Q}/\texttt{Esc} para salir,
\texttt{P} para pausar, \texttt{1}--\texttt{5} para conmutar cintas,
\texttt{R} para reiniciar, \texttt{S} para guardar el mosaico y
\texttt{T}/\texttt{H} para controlar la pose de prueba del robot cuando está
activado el modo IK manual.

\texttt{CintaTransportadora} mantiene su propio estado, velocidad y dirección.
\texttt{CamaraInspeccion} captura el render de la cámara MuJoCo, segmenta la
imagen en HSV, busca contornos externos y clasifica el mayor contorno válido
mediante circularidad y aproximación poligonal. Descarta contornos con área
inferior a 90 o superior a 26000 píxeles, así como dimensiones menores de 8
píxeles. Puede devolver círculo, triángulo, cuadrado, rectángulo o
desconocido.

La variable booleana \texttt{objeto\_detectado} indica si el procesamiento
actual encontró una forma. La planta detecta el flanco de una detección,
detiene la cinta asociada y la reanuda tras 0,35 segundos de simulación,
siempre que no exista una parada manual. Separar las paradas manuales del
bloqueo temporal de inspección evita que la reanudación automática anule una
orden manual.

\texttt{Horno} recibe piezas en una cola, mantiene los contadores y modifica su
temperatura. Parte de 20\,°C, con límites de 15\,°C y 250\,°C; su color lógico
es azul por debajo de 80\,°C, naranja entre 80\,°C y 180\,°C, y rojo desde
180\,°C. Cuando hay piezas en cola, la actualización aumenta la temperatura
0,01\,°C por paso y procesa una pieza al alcanzar 180\,°C; sin piezas, la
temperatura disminuye hacia el mínimo. El color se aplica a la geometría
\texttt{oven\_hot\_geom} mediante \texttt{model.geom\_rgba}.

\texttt{RobotPanda} envuelve las llamadas del controlador ya existente. El
código de configuración mantiene activado por defecto \texttt{ROBOT\_IK\_TEST\_MODE},
por lo que la ejecución inicial usa el control manual de la pose IK de prueba;
la rutina automática de recogida y colocación queda disponible en la rama de
ejecución cuando ese modo se desactiva.

\subsection{Fragmentos de código relevantes}

El siguiente fragmento resume el criterio de color implementado en el horno:

\begin{verbatim}
if self.temperatura < 80:
    self.color = "azul"
elif self.temperatura < 180:
    self.color = "naranja"
else:
    self.color = "rojo"
\end{verbatim}

La implementación completa se consulta en los módulos del proyecto y deberá
publicarse en el repositorio GitHub indicado en la versión final del informe.

% =========================================================
% 7. MANEJO DE ERRORES Y VALIDACIONES
% =========================================================

\section{Manejo de errores y validaciones}

\begin{itemize}
    \item Las clases validan dimensiones de posiciones y desplazamientos, dirección de la cinta, límites del horno e incrementos de temperatura; los valores fuera de rango producen \texttt{ValueError}.
    \item Se comprueba que la planta esté inicializada antes de actualizarla y que el robot esté vinculado al modelo antes de controlarlo; en caso contrario se genera \texttt{RuntimeError}.
    \item La inicialización comunica mediante errores explícitos la ausencia de la geometría del horno o de un actuador requerido.
    \item Si no se puede escribir una captura, se informa mediante \texttt{OSError}.
    \item Los errores no se ocultan con valores de sustitución: se propagan para que el problema pueda identificarse.
\end{itemize}

% =========================================================
% 8. PRUEBAS
% =========================================================

\section{Pruebas y validación}

No se encontraron pruebas automatizadas en los archivos disponibles. La tabla
siguiente es un plan de comprobación derivado de los requisitos; sus resultados
deben marcarse como realizados únicamente después de ejecutar y observar cada
caso en el entorno de entrega.

\begin{table}[H]
    \centering
    \caption{Plan de pruebas funcionales pendiente de ejecución y registro.}
    \label{tab:pruebas}
    \begin{tabularx}{\textwidth}{p{1.4cm}X X p{2.4cm}}
        \toprule
        \textbf{Prueba} & \textbf{Entrada / escenario} & \textbf{Resultado esperado} & \textbf{Estado} \\
        \midrule
        P01 & Iniciar \texttt{python main.py}. & Arrancan la simulación MuJoCo y la ventana OpenCV. & Pendiente \\
        P02 & Conmutar cintas con las teclas \texttt{1}--\texttt{5}. & Cada cinta cambia su estado sin conmutar las demás. & Pendiente \\
        P03 & Presentar piezas a las cámaras. & Se actualiza \texttt{objeto\_detectado} y se clasifica una forma visible. & Pendiente \\
        P04 & Mantener una detección y esperar su liberación. & La cinta se detiene durante la inspección y se reanuda salvo parada manual. & Pendiente \\
        P05 & Hacer llegar una pieza al horno y observar su evolución. & Cambian temperatura y color; al alcanzar el umbral se procesa una pieza. & Pendiente \\
        P06 & Usar \texttt{T} y \texttt{H} en modo IK de prueba. & El Panda recibe las poses mediante el controlador proporcionado. & Pendiente \\
        \bottomrule
    \end{tabularx}
\end{table}

% =========================================================
% 9. RESULTADOS
% =========================================================

\section{Resultados}

El proyecto integra las clases descritas con el modelo físico existente y
expone controles interactivos para las cintas, la pausa, el reinicio, la
captura de imágenes y el movimiento de prueba del robot. El comportamiento
descrito corresponde a la inspección del código; la ejecución completa de los
casos de prueba queda pendiente de registrar.

\begin{figure}[H]
    \centering
    \IfFileExists{capturas/mosaico.png}{%
        \includegraphics[width=0.92\textwidth]{capturas/mosaico.png}%
    }{%
        \fbox{%
            \parbox[c][4cm][c]{0.82\textwidth}{%
                \centering
                No hay una captura de ejecución en \texttt{capturas/}.\\
                Guardar el mosaico con la tecla \texttt{S} y añadirlo como\\
                \texttt{capturas/mosaico.png} antes de entregar el PDF.
            }%
        }%
    }
    \caption{Mosaico de cámaras y vista general de la planta.}
    \label{fig:resultado}
\end{figure}

% =========================================================
% 10. DISCUSIÓN
% =========================================================

\section{Discusión}

\begin{itemize}
    \item \textbf{Separación de responsabilidades:} las clases de dominio agrupan estado y operaciones propios, mientras la planta coordina la simulación y la interfaz.
    \item \textbf{Integración gradual:} se conservan las funciones físicas y el solucionador IK de partida, evitando reescribir componentes que no forman parte del objetivo de la práctica.
    \item \textbf{Acoplamiento:} \texttt{PlantaIndustrial} conoce el modelo MuJoCo y el módulo procedural auxiliar; esta dependencia facilita coordinar el ciclo actual, aunque limita la independencia del coordinador para pruebas unitarias.
    \item \textbf{Visión:} la clasificación por umbrales de HSV y geometría de contornos es sencilla y explicable, pero depende de las condiciones de iluminación y segmentación de la escena. La salida contempla formas geométricas, no una identificación semántica robusta de cada pieza.
    \item \textbf{Mejoras futuras:} añadir pruebas automatizadas, separar la política de decisión de la cámara, configurar los parámetros de visión por cámara y ampliar la coordinación del robot para la Práctica II.
\end{itemize}

% =========================================================
% 11. CONCLUSIONES
% =========================================================

\section{Conclusiones}

El proyecto estructura mediante objetos los componentes principales de una
planta simulada: cintas independientes, cámaras asociadas, piezas, horno,
coordinador y adaptador del robot. La cámara comunica su detección mediante
\texttt{objeto\_detectado}; el coordinador usa esa señal para detener
temporalmente la cinta, respetando las paradas manuales. El horno encapsula su
temperatura y comunica su estado térmico mediante el color de una geometría
MuJoCo.

La solución mantiene el simulador y el controlador IK existentes, en lugar de
reimplementar su física o cinemática. Esto proporciona una base modular para
extensiones posteriores, aunque persiste acoplamiento con las rutinas
procedurales originales. Para cerrar la entrega se deben ejecutar y registrar
las pruebas, incorporar una captura real, completar los datos personales y
añadir el enlace al repositorio.

% =========================================================
% 12. CONTRIBUCIÓN INDIVIDUAL
% =========================================================

\section{Contribución individual}

La información de los archivos del proyecto no permite identificar el nombre
del autor ni el reparto de tareas. Completar esta sección con la contribución
real de cada integrante antes de entregar. Si el trabajo es individual,
indicarlo expresamente y describir las tareas realizadas por el estudiante.

% =========================================================
% 13. USO DE INTELIGENCIA ARTIFICIAL
% =========================================================

\section{Declaración sobre el uso de herramientas de inteligencia artificial}

\begin{quote}
Se utilizó un asistente de inteligencia artificial como apoyo para redactar esta
memoria a partir del enunciado y de los archivos del proyecto. El estudiante
debe revisar el contenido, comprobar que describe fielmente su trabajo y
adaptarlo a las pruebas y contribuciones que haya realizado antes de entregar.
\end{quote}

% =========================================================
% 14. REFERENCIAS
% =========================================================

\section{Referencias}

\begin{thebibliography}{9}

\bibitem{enunciado}
Universidad de Burgos,
\textit{Práctica I: Refactorización de una Planta Industrial mediante
Programación Orientada a Objetos}, enunciado de la asignatura,
curso 2026--2027.

\bibitem{uml}
Object Management Group, \textit{Unified Modeling Language (UML)},
\url{https://www.omg.org/spec/UML/}.

\bibitem{mujoco}
Google DeepMind, \textit{MuJoCo Documentation},
\url{https://mujoco.readthedocs.io/}.

\bibitem{opencv}
OpenCV, \textit{OpenCV-Python Tutorials},
\url{https://docs.opencv.org/}.

\bibitem{python}
Python Software Foundation, \textit{Python Documentation},
\url{https://docs.python.org/3/}.

\end{thebibliography}

% =========================================================
% ANEXOS
% =========================================================

\newpage
\appendix

\section{Código fuente y repositorio}

El código debe publicarse en un repositorio GitHub accesible para la evaluación.
El enlace no se encuentra en los archivos disponibles y debe completarse antes
de entregar: \textbf{Enlace al repositorio: pendiente de añadir}.

Para instalar las dependencias desde la raíz del proyecto:
\begin{verbatim}
python -m pip install -r requirements.txt
\end{verbatim}

Para ejecutar la aplicación:
\begin{verbatim}
python main.py
\end{verbatim}

El modelo Franka se incluye en \texttt{vendor/mujoco\_menagerie}; el README
también documenta la variable \texttt{MUJOCO\_MENAGERIE\_PATH} para indicar
otra instalación.

\section{Material adicional}

La carpeta \texttt{capturas/} está prevista para las imágenes guardadas durante
la simulación. Incorporar aquí o en la sección de resultados la captura
solicitada por el enunciado y, si procede, material adicional de las pruebas.

\end{document}
